\chapter[RESULTADOS]{Resultados}
\label{cap:resultados}

Este capítulo apresenta os resultados de cada etapa do procedimento de
calibração descrito na Seção~\ref{sec:calibracao}: os diagnósticos da análise
de sensibilidade, as variantes estruturais avaliadas, a busca em grade e sua
validação, a configuração final e, por fim, a avaliação no conjunto de teste.
O capítulo termina com a discussão dos resultados.

% -----------------------------------------------------------------------------
\section{Análise de sensibilidade}
\label{sec:res_sensibilidade}

A análise OAT da Etapa C avaliou os 36 parâmetros numéricos da configuração
original. A investigação dos parâmetros de amplitude nula revelou três defeitos
estruturais:

\begin{enumerate}
  \item Os pesos da extração de prazo não tinham efeito porque a biblioteca
    \textit{dateparser}, que recebe a frase inteira (Seção~\ref{subsec:prazo}), não
    interpreta datas dentro de frases: ela retornou nulo para todas as 976
    frases candidatas da calibração. Na prática, todas as datas vinham do
    \textit{fallback} por expressão regular, e a pontuação de frases era
    inoperante. Por isso, a acurácia da data na configuração original foi
    nula.
  \item A regra de contagem de \textit{links} do filtro de listagem procurava
    \textit{links} no formato Markdown, mas o texto analisado vem de HTML
    convertido, que não tem esse formato. A regra nunca era acionada.
  \item O limiar secundário de relevância (2) era sempre menor que o limiar
    principal (3), aplicado antes dele, e portanto nunca decidia.
\end{enumerate}

Esses defeitos não eram visíveis nas métricas agregadas: o sistema produzia
resultados plausíveis, e a falha só apareceu porque a análise expôs
parâmetros que não alteravam nenhum resultado.

% -----------------------------------------------------------------------------
\section{Variantes estruturais}
\label{sec:res_variantes}

A Tabela~\ref{tab:variantes} resume as variantes avaliadas na Etapa D e as
decisões tomadas pelo critério lexicográfico.

\begin{table}[htbp]
\centering
\caption{Variantes estruturais avaliadas no conjunto de calibração}
\label{tab:variantes}
\small
\begin{tabularx}{\textwidth}{>{\raggedright\arraybackslash}p{3.6cm} X >{\raggedright\arraybackslash}p{3.8cm}}
\toprule
Variante & Hipótese & Decisão \\
\midrule
Datas dentro da frase
  & Reconhecer datas por expressão regular dentro de cada frase candidata, no
    lugar da interpretação da frase inteira
  & Aceita \\
Prazo pela data de maior pontuação
  & Decidir se o edital está aberto pela data mais bem pontuada, e não pela
    melhor data futura
  & Aceita \\
Data por segmento de palavra-chave
  & Em frases com várias datas, usar a última data do primeiro segmento
    iniciado por um termo de prazo
  & Rejeitada: a acurácia da data caiu de 9/19 para 7/19 \\
Data por âncora de limite
  & Em frases com várias datas, usar a primeira data após ``limite'',
    ``encerramento'' ou similar
  & Aceita \\
Listagem sem \texttt{/editais/} e \texttt{/chamadas/}
  & Esses segmentos de URL identificam páginas de editais individuais, não de
    listagens
  & Aceita: editais perdidos de 7 para 4 \\
Listagem sem as palavras ``editais'' e ``chamadas'' na URL
  & Mesma hipótese, estendida às palavras de URL
  & Rejeitada: os itens errados aceitos passaram de 2 para 4 \\
Vocabulário temático de relevância
  & Acrescentar termos como ``fomento'', ``subvenção'' e ``pesquisa''
  & Rejeitada: a $F_2$ caiu \\
Remoção do limiar secundário de relevância
  & O limiar nunca decide
  & Aceita: métricas idênticas, código mais simples \\
Remoção da regra de \textit{links}
  & A regra nunca é acionada
  & Aceita: métricas idênticas \\
\bottomrule
\end{tabularx}
\fonte{Elaborado pelo autor.}
\end{table}

As três variantes aceitas na extração de prazo atuam em conjunto. Isoladamente,
o reconhecimento de datas dentro da frase elevou a acurácia da data de 0/19
para 9/19, mas aumentou os itens errados aceitos de 2 para 5. A análise desses
itens mostrou que, em três dos quatro editais encerrados aceitos por engano, a
data mais bem pontuada era a data correta, já vencida, e o item era aceito
porque a regra original escolhia a melhor data \emph{futura}, em geral uma data
de etapa posterior, como a divulgação de resultados. A variante de prazo pela
data de maior pontuação corrige esse comportamento: o item é considerado
encerrado quando a data mais bem pontuada já passou.

A regra da âncora de limite responde a uma limitação da extração de texto. A
conversão do HTML colapsa todos os espaços em branco, inclusive as quebras de
linha. Com isso, tabelas de cronograma e menus de navegação tornam-se
``frases'' muito longas, que reúnem várias datas e termos de etapas diferentes.
Nessas frases, para cada ocorrência de ``limite'', ``encerramento'',
``término'', ``retificado'' ou ``prorrogado'', toma-se a primeira data que a
sucede, desde que não haja um termo negativo entre as duas, e prevalece a
última âncora que tem data. Na ausência de âncoras, adota-se a última data da
frase, que corresponde ao fim do período em construções como ``de X a Y''. Com
as três variantes, a acurácia da data na calibração chegou a 10/19, com 2 itens
errados aceitos, o mesmo número da configuração original.

Antes de remover a regra de \textit{links}, verificou-se se ela poderia ser
corrigida, contando os \textit{links} diretamente no HTML bruto. A contagem não
discriminou as classes: a mediana foi de 90 \textit{links} por página nas
listagens e de 78 nas demais, pois os menus de navegação dos portais
institucionais inserem centenas de \textit{links} em qualquer página. Optou-se
pela remoção, em vez de calibrar um limiar sobre um sinal fraco.

% -----------------------------------------------------------------------------
\section{Busca em grade e validação}
\label{sec:res_grade}

Na repetição da análise OAT sobre a base com as variantes aceitas, todos os
valores vigentes da extração de prazo estavam no patamar ótimo. No filtro de
listagem, vários parâmetros melhoravam o resultado em um item cada, o que
indicava possíveis efeitos combinados e justificou a busca em grade.

A combinação escolhida pela grade alterou quatro parâmetros em relação à base:
$w_p = 0$, $d_{\mathrm{high}} = 20$, $w_e = 2$ e $\theta_L = 6$. Com ela, o
número de editais perdidos na calibração caiu de 4, após as variantes, para 1,
sem aumento dos itens errados aceitos.

Na validação cruzada, a configuração escolhida em cada partição perdeu 32
editais nos itens retidos, somadas as 50 partições, contra 40 da base, com 22
itens errados aceitos contra 20. A mesma combinação foi selecionada em 32 das
50 partições, e as quatro combinações mais frequentes, que somam 45 partições,
coincidem em quatro dos seis parâmetros. A escolha foi considerada
suficientemente estável e adotada. O pequeno aumento de itens errados aceitos
fora da amostra de ajuste indica, entretanto, que parte do ganho observado na
calibração pode não se repetir em dados novos.

No \textit{bootstrap} sobre o conjunto de calibração, o intervalo de confiança
de 95\% da diferença de editais perdidos entre a configuração calibrada e a
original foi de $-11$ a $-2$, e o da diferença de itens errados aceitos, de
$-4$ a $4$.

% -----------------------------------------------------------------------------
\section{Configuração final}
\label{sec:configuracao}

A calibração alterou quatro parâmetros e a lista de padrões de URL do filtro de
listagem, ativou três variantes da extração de prazo e removeu dois mecanismos
sem efeito. Os demais parâmetros permaneceram nos valores originais. A
Tabela~\ref{tab:configuracao} resume as diferenças.

\begin{table}[htbp]
\centering
\caption{Parâmetros alterados pela calibração}
\label{tab:configuracao}
\small
\begin{tabularx}{\textwidth}{l >{\raggedright\arraybackslash}X >{\raggedright\arraybackslash}p{2.6cm} >{\raggedright\arraybackslash}p{2.6cm} l}
\toprule
Componente & Parâmetro & Original & Calibrado & Origem \\
\midrule
Listagem   & Padrões de URL & com \texttt{/editais/} e \texttt{/chamadas/} & sem \texttt{/editais/} e \texttt{/chamadas/} & Variante \\
Listagem   & Peso do padrão de URL ($w_p$) & 3 & 0 & Grade \\
Listagem   & Segunda faixa de datas ($d_{\mathrm{high}}$) & 10 & 20 & Grade \\
Listagem   & Peso do plural ``editais'' ($w_e$) & 1 & 2 & Grade \\
Listagem   & Limiar ($\theta_L$) & 5 & 6 & Grade \\
Listagem   & Regra de contagem de \textit{links} & ativa (3 parâmetros) & removida & Etapa C \\
Relevância & Limiar secundário & 2 & removido & Etapa C \\
Prazo      & Datas dentro da frase & desligada & ligada & Variante \\
Prazo      & Prazo pela data de maior pontuação & desligada & ligada & Variante \\
Prazo      & Data por âncora de limite & desligada & ligada & Variante \\
\bottomrule
\end{tabularx}
\fonte{Elaborado pelo autor.}
\end{table}

Com o peso $w_p$ igual a 0, os padrões de URL restantes deixaram de influenciar
a pontuação de listagem. O parâmetro foi mantido na configuração para preservar
a compatibilidade com experimentos futuros. Na configuração final, a biblioteca
\textit{dateparser} passou a ser usada apenas no \textit{fallback}.

% -----------------------------------------------------------------------------
\section{Avaliação no conjunto de teste}
\label{sec:res_teste}

A Tabela~\ref{tab:resultados} compara a configuração original e a calibrada
nos dois conjuntos. As colunas de calibração mostram o ajuste obtido nos dados
usados para calibrar; as de teste, o desempenho em dados não vistos.

\begin{table}[htbp]
\centering
\caption{Desempenho das configurações original e calibrada}
\label{tab:resultados}
\small
\begin{tabular}{l c c c c}
\toprule
 & \multicolumn{2}{c}{Calibração (47 itens)} & \multicolumn{2}{c}{Teste (26 itens)} \\
\cmidrule(lr){2-3} \cmidrule(lr){4-5}
Medida & Original & Calibrada & Original & Calibrada \\
\midrule
Editais perdidos                 & 7/20  & 1/20  & 5/14  & 1/14 \\
Acurácia da data                 & 0/19  & 10/19 & 0/12  & 4/12 \\
Itens errados aceitos            & 2     & 2     & 2     & 2    \\
Aceitos com a data certa         & 0/3   & 0/2   & 0/2   & 1/3  \\
Resultado final ($VP$/positivos) & 1/1   & 0/1   & 0/2   & 1/2  \\
\midrule
Listagem: precisão               & 0,53  & 0,75  & 0,29  & 0,50 \\
Listagem: revocação              & 0,56  & 0,33  & 0,40  & 0,20 \\
Listagem: $F_{0{,}5}$            & 0,53  & 0,60  & 0,30  & 0,38 \\
Relevância: $F_2$                & 0,90  & 0,90  & 0,67  & 0,67 \\
Prazo aberto: $F_1$              & 0,29  & 0,00  & 0,22  & 0,67 \\
\bottomrule
\end{tabular}
\fonte{Elaborado pelo autor.}
\end{table}

No conjunto de teste, a configuração calibrada perdeu 1 dos 14 editais, contra
5 da original. O intervalo de confiança de 95\% da diferença, estimado por
\textit{bootstrap}, foi de $-8$ a $-1$. A acurácia da data passou de 0/12 para
4/12, e a medida $F_1$ do prazo aberto, de 0,22 para 0,67. O número de itens
errados aceitos permaneceu em 2, com intervalo de confiança da diferença de
$-4$ a $4$. Na avaliação por etapa, a configuração calibrada aceitou três
itens, dos quais um era um edital aberto e relevante, entre os dois existentes
no conjunto; a configuração original aceitou dois itens, nenhum deles correto.
Na execução em cadeia, a deduplicação por conteúdo descartou um dos itens
errados, por similaridade de 0,86 com outro edital da mesma agência, e o
sistema entregou dois itens, um deles correto. O filtro de relevância não foi
alterado e manteve as mesmas métricas.

A distribuição das 2.000 diferenças reamostradas
(Equação~\ref{eq:bootstrap_reamostra}) explica a diferença entre os dois
intervalos. Nos editais perdidos, a diferença foi negativa em 98,7\% das reamostras e nula nas demais; nenhuma
reamostra indicou piora, porque, item a item, a configuração calibrada não
perdeu nenhum edital que a original não perdesse. Nos itens errados aceitos, a
distribuição é aproximadamente simétrica em torno de zero: a diferença foi
negativa em 39,9\% das reamostras, nula em 21,2\% e positiva em 39,0\%. O
total igual de itens errados resulta de uma troca, pois as duas configurações
aceitaram indevidamente itens diferentes. Os dados por item que produzem as
duas distribuições estão no Apêndice~\ref{ap:bootstrap}.

% -----------------------------------------------------------------------------
\section{Discussão}
\label{sec:discussao}

A redução de editais perdidos foi o resultado mais consistente da calibração.
Ela ocorreu nos dois conjuntos, em proporção semelhante, e em ambos o intervalo
de confiança da diferença exclui o zero. A principal causa foi a remoção dos
segmentos \texttt{/editais/} e \texttt{/chamadas/} dos padrões de listagem, que
classificavam como listagem justamente as páginas de editais individuais.

O ganho de editais teve como contrapartida a queda da revocação do filtro de
listagem, de 0,56 para 0,33 na calibração e de 0,40 para 0,20 no teste: mais
listagens passaram a atravessar o filtro. Essa troca é coerente com o critério
lexicográfico adotado, que prioriza não descartar editais, e com o desenho em
cadeia, no qual uma listagem aprovada ainda pode ser barrada pelos filtros
seguintes, enquanto um edital descartado é perdido. A precisão do filtro
aumentou nos dois conjuntos.

A acurácia da data foi a métrica com maior diferença entre calibração e teste:
10/19 contra 4/12. Parte da diferença decorre do tamanho das amostras, mas ela
também é compatível com sobreajuste da regra da âncora de limite, concebida a
partir dos erros observados nas páginas de calibração. Ainda assim, o ganho em
relação à configuração original, que não acertava nenhuma data, se manteve no
teste.

O número de itens errados aceitos não mudou em nenhum dos conjuntos. A
calibração, portanto, não melhorou a precisão do resultado final. Na
configuração calibrada, três dos quatro itens errados aceitos são listagens
que atravessaram o filtro de listagem e que, por conterem vocabulário de
fomento e datas futuras, também foram aprovadas pelos filtros de relevância e
de prazo; o quarto é um edital encerrado ao qual foi atribuída uma data futura
incorreta. Esse padrão é consequência direta da troca descrita acima e da
ausência de uma etapa que decida se a página é um edital
(Seção~\ref{sec:arquitetura}).

As métricas do resultado final devem ser lidas com cautela. Com um edital aberto
na calibração e dois no teste, um único item altera o resultado de 0 para 100\%
ou de 0 para 50\%. A queda de 1/1 para 0/1 na calibração, por exemplo, reflete
um acerto da configuração original obtido por coincidência: o único edital
aberto era aceito com uma data errada, e a extração calibrada, que escolhe
outra data igualmente errada, porém já vencida, passou a considerá-lo
encerrado. Pelo mesmo motivo, a
$F_1$ do prazo aberto caiu a zero na calibração e subiu no teste. Por isso, a
avaliação do trabalho apoia-se principalmente nas métricas por etapa e nas
contagens de editais perdidos e de itens errados aceitos.
